

%% !TEX root = ../main.tex
\section{Conclusão}

O objeto principal desta análise foi o \texttt{educado-web}, conforme
delimitado no escopo do trabalho; o \texttt{educado-api} foi examinado
apenas na medida em que o funcionamento do Web depende dela, sobretudo
para os fluxos que exigem testes de integração entre a interface e o
backend. Sob esse recorte, o achado central é que o \texttt{educado-web}
possui uma proteção mínima por testes automatizados: apenas um teste de
fumaça, cuja função é validar a infraestrutura de testes em si, e não os
comportamentos da aplicação~\cite{secao5}. Esse teste foi executado com
sucesso, o que confirma que a infraestrutura (Vitest, jsdom) está
corretamente configurada e integrada ao CI, mas isso não equivale a uma
suíte que proteja o produto.

A API, examinada como dependência do Web, mostrou-se consideravelmente
mais madura: 31 conjuntos e 506 testes aprovados, com cobertura de
79,36\% de instruções, 69,26\% de branches, 78\% de funções e 78,9\% de
linhas no escopo definido pelo \texttt{jest.config.ts}~\cite{secao5}.
Esse dado é relevante para o trabalho não porque a API seja o foco da
análise, mas porque ele delimita o que já está coberto do lado do
backend e, por consequência, o que falta ser verificado do lado do Web, em especial a integração entre os dois: o consumo da API pela
interface, o tratamento de respostas de erro, autenticação via JWT
armazenado no \texttt{localStorage} e a reação da aplicação a falhas de
rede.

Esse veredito, no entanto, cobre apenas a verificação automatizada de
código, e o próprio relatório identifica que essa é só parte do que o
Educado verifica. Existe uma frente de campo, ancorada em observações
diretas com catadores e registrada em \texttt{psp5} e no
\texttt{dashboard}~\cite{psp5,dashboard}, que avalia usabilidade e
aceitação da plataforma com usuários reais. Essa frente não é considerada
na avaliação acima porque seu foco está no aplicativo e na experiência de
conteúdo do estudante, não na interface administrativa e de criação de
cursos que compõe o escopo deste trabalho; ainda assim, sua existência
mostra que a lacuna de testes automatizados no Web não é a mesma coisa
que uma lacuna de verificação do produto como um todo, é uma lacuna
específica do tipo de verificação que este trabalho se propôs a analisar.

O ponto forte identificado que se aplica diretamente ao escopo do Web é o
processo de contribuição: o workflow de CI do \texttt{educado-web} já
executa os testes automaticamente em pull requests, e o modelo de pull
request do repositório solicita procedimentos de teste e evidências
visuais~\cite{secao5}. Ou seja, a estrutura para sustentar uma suíte
maior já existe; falta populá-la.

A principal limitação, portanto, não é a ausência de testes no projeto
como um todo, a API demonstra que a equipe sabe construir uma suíte
consistente, e a frente de campo mostra que a verificação com usuários
reais também acontece,, mas a concentração quase total da verificação
automatizada de código fora do escopo deste trabalho. Para o
\texttt{educado-web} especificamente, isso significa que funcionalidades
centrais para os usuários administradores e criadores de conteúdo,
autenticação, criação e edição de cursos, roteamento, internacionalização
e comunicação com a API, podem sofrer regressões sem que o CI as
detecte~\cite{secao5}.

As melhorias mais importantes para o escopo deste trabalho são, em
ordem de prioridade: (i) ampliar a suíte do \texttt{educado-web} com
testes unitários e de componente (Vitest, jsdom, Testing Library),
priorizando autenticação, formulários, permissões e o cliente HTTP; (ii)
adicionar testes de integração entre Web e API para os fluxos que
efetivamente dependem do backend, exatamente o ponto em que o escopo do
trabalho autoriza o uso da API,, cobrindo cenários de erro e expiração
de sessão; e (iii) complementar com uma suíte enxuta de testes
\textit{end-to-end} para os fluxos mais críticos de login, criação e
publicação de curso~\cite{secao5}.

Em síntese, a equipe avalia a estratégia de \emph{testes automatizados de
código} do \texttt{educado-web} como \textbf{insuficiente diante do
escopo do produto}, ainda que apoiada por uma infraestrutura de CI já
funcional e por uma API bem testada que pode servir de referência de
maturidade a ser alcançada. Esse veredito é específico à camada de
código: o projeto complementa a verificação com uma frente de campo
voltada a usuários reais, que este trabalho não teve como escopo avaliar
em profundidade. O uso pontual da API neste trabalho reforça o
diagnóstico sobre a camada de código: a maior lacuna não está em nenhum
dos dois lados isoladamente, mas na fronteira entre eles, exatamente o
tipo de teste que o escopo do trabalho previu como necessário e que
ainda não existe.